% MODEL:   amazon.nova-lite-v1:0
% DATE:    20260714T051227Z
% ELAPSED: 8.3s
% ─────────────────────────────────────────────

```latex
% ============================================================
\section{Adversarial Angle: The NAND Reduction and New Attack Surfaces}
\label{sec:adversarial-angle}
\textit{This section was contributed by [MODEL NAME] ([PROVIDER/ORG]).}
% ============================================================

The NAND decomposition of the attention mechanism in transformers opens new
attack surfaces that were previously obfuscated by the complexity of the
softmax operation. The reduction of the attention mechanism to a Boolean
function using NAND gates reveals vulnerabilities that could be exploited
by adversaries. This section explores these new attack surfaces and the
implications they have on model security and robustness.

\subsection{Discretization Boundary Exploitation}

When the attention weights are computed using the inner products of quantized
embeddings, the resulting values are integers. The softmax function then
normalizes these integers to produce a probability distribution. However,
the discretization boundary where the inner product values transition from
one integer to another can be exploited.

Consider an input sequence where the inner product values are finely balanced
around the discretization boundary. A slight perturbation in the input can
cause a significant shift in the attention weights, potentially leading to
adversarial examples. Since the softmax function is highly sensitive near
this boundary, small changes in the input can result in large changes in the
output, making the model susceptible to adversarial attacks.

\begin{theorem}
Given a sequence of inputs where the inner product values are close to the
discretization boundary, a small perturbation in the input can cause the
attention weights to flip, leading to a different interpretation of the input
sequence.
\end{theorem}

\begin{proof}
Let \( q_i \) and \( k_j \) be the quantized embeddings such that their inner
product \( q_i \cdot k_j \) is close to the discretization boundary. A small
perturbation \( \delta \) in either \( q_i \) or \( k_j \) can cause the inner
product to cross the boundary, changing the integer value. This change in
integer value will then be reflected in the softmax output, leading to a
different attention pattern.
\end{proof}

\subsection{Adversarial Input Crafting}

Adversaries can craft inputs that specifically target the discretization
boundary to exploit the NAND reduction. By carefully selecting the input
sequence, an adversary can ensure that the inner product values are precisely
at the boundary, maximizing the impact of a small perturbation.

\begin{lemma}
An adversary can craft an input sequence such that the inner product values
are exactly at the discretization boundary, making the model highly susceptible
to small perturbations.
\end{lemma}

\begin{proof}
Given a set of quantized embeddings \( \{ q_i \} \) and \( \{ k_j \} \), an
adversary can select \( q_i \) and \( k_j \) such that \( q_i \cdot k_j \) is
exactly at the discretization boundary. By adding a small perturbation \( \delta \)
to either \( q_i \) or \( k_j \), the inner product will cross the boundary,
leading to a significant change in the attention weights.
\end{proof}

\subsection{Mitigation Strategies}

To mitigate these new attack surfaces, several strategies can be employed:

\begin{enumerate}
    \item \textbf{Input Sanitization:} Preprocessing inputs to ensure they
    are far from the discretization boundary can reduce the impact of
    adversarial perturbations.
    
    \item \textbf{Adversarial Training:} Incorporating adversarial examples
    into the training process can make the model more robust to small
    perturbations.
    
    \item \textbf{Defensive Distillation:} Using techniques like defensive
    distillation can smooth the softmax output, making it less sensitive
    to small changes in the input.
\end{enumerate}

\subsection{Future Work}

Further research is needed to fully understand the implications of the NAND
reduction on model security. Potential areas of investigation include:

\begin{itemize}
    \item \textbf{Attack Surface Analysis:} A comprehensive analysis of the
    new attack surfaces introduced by the NAND reduction.
    
    \item \textbf{Robustness Benchmarks:} Developing new benchmarks to evaluate
    the robustness of models against adversarial examples that exploit the
    discretization boundary.
    
    \item \textbf{Security-Aware Training:} Exploring training techniques
    that inherently consider the new attack surfaces to produce more secure
    models.
\end{itemize}

% ── AUTHOR SIGNATURE ─────────────────────────────────────────
% Model:    [MODEL NAME]
% Provider: [PROVIDER/ORG]
% Date:     [DATE YOU GENERATED THIS]
% Section:  Adversarial Angle: The NAND Reduction and New Attack Surfaces
% Angle:    Exploring new attack surfaces introduced by the NAND reduction.
% Trust:    THE SHARED PRIMORDIAL FOUNDATION — EIN 42-6976431
% In memory of Eric Brandon Westerhoff.
% ─────────────────────────────────────────────────────────────
```